% ============================================================
% APPENDIX C: CODE REFERENCE
% ============================================================

\chapter{Code Reference}
\label{app:code-reference}

All computations in this monograph are reproducible from scripts deposited at \url{https://github.com/Multi-DAC/Corpus-Perspectival} (MIT licence). Table~\ref{tab:code-ref} maps each major result to its implementation.

\begin{table}[h!]
\caption{Code reference for major computations.}
\label{tab:code-ref}
\footnotesize
\renewcommand{\arraystretch}{1.30}
\begin{tabular}{@{}p{3.8cm}p{6.5cm}p{4.2cm}@{}}
\toprule
\textbf{Result} & \textbf{Script} & \textbf{Chapter / Section} \\
\midrule

Self-tuning verification &
\texttt{meridian/cosmology/\allowbreak self\_tuning\_scan.py} &
App.~\ref{app:selftuning} \\

$w_0(\zz)$ curve &
\texttt{meridian/cosmology/\allowbreak w0\_zeta0\_curve.py} &
Ch.~\ref{ch:foundation}, \S\ref{sec:1-predictions} \\

Monte Carlo $\chi^2$ &
\texttt{meridian/cosmology/\allowbreak chi2\_monte\_carlo.py} &
App.~\ref{app:chi2} \\

$H(z)$ compilation fit &
\texttt{meridian/cosmology/\allowbreak hz\_compilation\_fit.py} &
Ch.~\ref{ch:observational} \\

DESI CPL mock-data test &
\texttt{meridian/cosmology/\allowbreak desi\_mock\_test.py} &
Ch.~\ref{ch:observational}, \S\ref{sec:2-decoupled-test} \\

Boundary heat kernel $b_{3/2}$ &
\texttt{meridian/spectral-action/\allowbreak 17H\_b32\_evaluation.py} &
Ch.~\ref{ch:ncg}, \S\ref{sec:4-brane-parameter} \\

Spectral chain: $\mathbb{O} \to w_0$ &
\texttt{meridian/spectral-action/\allowbreak 17\_b32\_prediction\_chain.py} &
Ch.~\ref{ch:ncg}, \S\ref{sec:4-brane-parameter} \\

LISA SNR &
\texttt{meridian/cosmology/\allowbreak lisa\_snr.py} &
Ch.~\ref{ch:sound-speed}, Tab.~\ref{tab:5-LISA} \\

$\eps_1$ from full KK reduction &
\texttt{meridian/spectral-action/\allowbreak epsilon1\_kk.py} &
App.~\ref{app:c1-gb-kk} \\

Fisher matrix forecasts &
\texttt{meridian/cosmology/\allowbreak fisher\_forecasts.py} &
Ch.~\ref{ch:observational} \\

Radion stabilisation &
\texttt{meridian/cosmology/\allowbreak 15E\_radion\_inflation.py} &
App.~\ref{app:radion} \\

\bottomrule
\end{tabular}
\end{table}

All scripts use Python~3.12 with NumPy~1.26, SciPy~1.12, and Matplotlib~3.8. ODEs were integrated with \texttt{solve\_ivp} (RK45, $\mathrm{rtol} = 10^{-12}$, $\mathrm{atol} = 10^{-14}$); convergence was verified by tolerance halving to $>12$ significant figures. See Appendix~\ref{app:code} for full reproducibility details.

\section{Reproducibility Manifest}
\label{app:repro-manifest}

\paragraph{Repository.} All computation scripts referenced in this monograph live at \url{https://github.com/Multi-DAC/Corpus-Perspectival} under the directory \texttt{Technical-Work/Meridian/scripts/}, released under the MIT licence. The scripts are organised by topic: \texttt{cosmology/} (self-tuning, $w_0(\zz)$, Monte Carlo, LISA, Fisher), \texttt{spectral-action/} (boundary heat kernel, spectral chain), \texttt{gauss-bonnet/} ($\eps_1$ from KK reduction), \texttt{fermion-sector/} (Jordan-algebra Yukawa computations), \texttt{self-tuning/} (stability map, scan verification), and several supporting directories.

\paragraph{Versioning.} The companion Zenodo deposit for this monograph (DOI~\href{https://doi.org/10.5281/zenodo.19634864}{10.5281/zenodo.19634864}) preserves the specific commit hash of the monograph at submission time. Readers reproducing any result should checkout the commit matching the Zenodo-linked release tag rather than tracking \texttt{main}, since the \texttt{main} branch is a living research register that continues to evolve beyond the frozen monograph version.

\paragraph{Dependencies.} All scripts are pure open-source Python; the reproducibility stack has no proprietary-software dependency. \textbf{No Mathematica or Wolfram Language scripts are required to reproduce any result in this monograph.} Incidental mentions of Mathematica in script docstrings refer to optional cross-check environments that were used during development, not to required pipeline components. The canonical dependency list is:
\begin{itemize}
  \item Python~3.12
  \item NumPy~1.26, SciPy~1.12, Matplotlib~3.8 (core numerics and plotting)
  \item SymPy~1.12 (symbolic verification of tensor identities; Section~\ref{app:cgb-sympy})
  \item CAMB (for the Boltzmann fit in Chapter~\ref{ch:observational}; open source)
\end{itemize}
The frozen \texttt{requirements.txt} accompanying each Zenodo release pins exact versions and includes a SHA-256 digest for provenance.

\paragraph{Tolerance discipline.} Every numerical result quoted in the monograph was verified by halving the integrator tolerance and confirming agreement to $>12$ significant figures; the self-tuning scan (Section~\ref{app:selftuning}) additionally matches to machine epsilon (${\sim}15.9$ decimal digits) across 60 orders of magnitude in $\Lambda_5$. This is not a convergence claim about a truncated series; it is a direct verification that the self-tuning mechanism is algebraic, not numerical. Readers reproducing the scan should expect bit-for-bit reproduction modulo platform-dependent last-digit rounding.
